\section{A canonical log-gamma and both Binet formulas}
\label{sec:binet}

\subsection{The right-half-plane logarithm}

For $z=x+iy$ with $x\ge a>0$, define
\[
 \operatorname{Log}z
 =\tfrac12\log(x^2+y^2)
    +i\int_0^{y/x}\frac{du}{1+u^2}.
\]
This is a primitive recursive name on each certified bounded rectangle.
Finite differentiation gives its complex derivative $1/z$.
For fixed $x>0$, differentiate
$\exp(\operatorname{Log}(x+iy))/(x+iy)$ in the real variable $y$:
the derivative is zero and its value at $y=0$ is one. The finite
fundamental theorem therefore gives $\exp(\operatorname{Log}z)=z$.
Positive scaling gives
$\operatorname{Log}(cz)=\log c+\operatorname{Log}z$ for $c>0$.
Define
\[
 \mathcal L_N(z)=\log(N!)+z\log N-
                        \sum_{k=0}^N\operatorname{Log}(z+k).
\]
For $N\ge2R$, where $|z|\le R$, its tail is bounded by
$(R^2+R)/N$: rewrite it as
\[
 -\operatorname{Log}z-\gamma_Nz
 -\sum_{k=1}^N\bigl(\operatorname{Log}(1+z/k)-z/k\bigr)
\]
and apply (C2)--(C6). On $|w|\le1/2$ the equality
$\operatorname{Log}(1+w)=w+L(w)$ is explicit: both have derivative
$1/(1+w)$ and value zero at $w=0$, so integrate their difference
along the segment $u\mapsto uw$. Denote the resulting effective limit by
$\mathcal L(z)$. Finite exponentiation and differentiation show
\begin{equation}
 e^{\mathcal L(z)}=\Gamma(z),\qquad
 \mathcal L'(z)=\psi(z),\qquad
 \mathcal L(z+1)-\mathcal L(z)=\operatorname{Log}z.
 \label{eq:canonical-loggamma}
\end{equation}
For the last equality the finite difference is
$\operatorname{Log}z+\log N-\operatorname{Log}(z+N+1)$,
whose last two terms tend explicitly to zero. On the positive axis
$\mathcal L(x)=\log\Gamma(x)$. In general $\mathcal L(z)$ need not be
the principal logarithm of the value $\Gamma(z)$; the normalization is
part of the definition.

\subsection{A positive kernel from reflection}

The differentiated reflection formula at $1/2$ and (C33) give
$3\zeta(2)=\pi^2/2$, hence $\zeta(2)=\pi^2/6$.
Subtract the digamma series (C27) at $z$ and $-z$, and use its recurrence
and reflection formula. Finite reindexing gives
\[
 \pi\cot(\pi z)=\frac1z-
                     2z\sum_{n=1}^{\infty}\frac1{n^2-z^2}.
\]
All derivative and series tails have already been bounded. Substituting
$z=it/(2\pi)$, with $t>0$, proves
\begin{equation}
 k(t):=\frac1t\left(\frac12-\frac1t+\frac1{e^t-1}\right)
       =2\sum_{n=1}^{\infty}\frac1{t^2+4\pi^2n^2}.
 \label{eq:binet-kernel}
\end{equation}
Use the series on the right to define $k(0)=1/12$ as well. It gives
\begin{equation}
 0<k(t)\le\frac1{12},\qquad
 \left|k(t)-2\sum_{n=1}^{N}\frac1{t^2+4\pi^2n^2}\right|
 \le\frac1{2\pi^2N}\le\frac1{8N}\quad(t\ge0).
 \label{eq:binet-kernel-tail}
\end{equation}
For strict positivity the first summand already supplies a rational
lower bound on every bounded interval. Thus the apparent cancellation
at $t=0$ is handled by a regular series rather than subtraction of
singular expressions.

Define
\[
 J(z)=\int_0^\infty e^{-zt}k(t)\dd t\qquad(\Rea z\ge a>0).
\]
The tail after $T$ is at most $e^{-aT}/(12a)$. On $[0,T]$ approximate
the kernel by $N$ summands; its integral error is at most $T/(8N)$.
The derivative of this finite integrand has a rational bound in terms
of $a,R,T$, and therefore a prescribed rational mesh completes the
algorithm. A sufficient coarse choice is a mesh Lipschitz bound
$R+T+1$, since the differentiated kernel series is bounded there by
$T$. Differentiation in $z$ adds powers of $t$; their tails are covered
by the monomial estimates used for Euler's integral.

Put $C=\tfrac12\log(2\pi)$ and
\[
 \mathcal R(z)=\mathcal L(z)-(z-\tfrac12)\operatorname{Log}z+z-C.
\]
Using \eqref{eq:binet-kernel} after multiplying by $t^2$, the finite
geometric sum for $(e^t-1)^{-1}$, and
$\int_0^\infty te^{-(z+n)t}\dd t=(z+n)^{-2}$, we get
\begin{align*}
 J''(z)
 &=\frac1{2z^2}-\frac1z+\sum_{n=1}^{\infty}\frac1{(z+n)^2}\\
 &=\psi'(z)-\frac1z-\frac1{2z^2}=\mathcal R''(z).
\end{align*}
The series tail is at most $1/(a+N)$, so the equality is justified
before passing to the infinite notation.

It remains to fix the possible affine difference; merely comparing
second derivatives would be insufficient. Rectangle-sum comparison for
$t\mapsto(w+t)^{-1}$ in the right half-plane gives
\begin{equation}
 |\psi(w)-\operatorname{Log}w|\le\frac1{\Rea w}.
 \label{eq:complex-psi-bound}
\end{equation}
Indeed the sum of the absolute derivative errors over $[0,N]$ is at
most $1/\Rea w-1/(\Rea w+N)$, the final endpoint contributes at most
$1/(\Rea w+N)$, and
$\log N-\operatorname{Log}(w+N)$ has an explicit vanishing bound.
Consequently $|\mathcal R'(w)|\le3/(2\Rea w)$. Integrate vertically
from $N+\Rea z$ to $N+z$ and use the proved real Stirling estimate:
\[
 |\mathcal R(N+z)|\le
 \frac{1/12+(3/2)|\Ima z|}{N+\Rea z},\qquad
 |J(N+z)|\le\frac1{12(N+\Rea z)}.
\]
The function $D=\mathcal R-J$ has zero second derivative, so finite
fundamental-theorem integration on a horizontal segment yields
\[
 D(z)=2D(z+N)-D(z+2N).
\]
The displayed bounds make the right side arbitrarily small with a
primitive recursive choice of $N$. Hence $D(z)=0$.

\begin{theorem}[Binet's first formula]
On every certified bounded region in $\Rea z>0$,
\[
 \boxed{\displaystyle
 \mathcal L(z)=(z-\tfrac12)\operatorname{Log}z-z
 +\tfrac12\log(2\pi)
 +\int_0^\infty
 \left(\frac12-\frac1t+\frac1{e^t-1}\right)
                 \frac{e^{-zt}}t\dd t.}
\]
At zero the integrand is interpreted using the regular kernel
\eqref{eq:binet-kernel}; all convergence estimates above are uniform
on the certified region.
\end{theorem}

\subsection{Binet's second formula without a contour deformation}

For real $t\ge0$ and $\Rea z>0$, specify the branch
\[
 \operatorname{atan}(t/z):=
                 \int_0^t\frac{z}{z^2+u^2}\dd u.
\]
The denominator factors $z+iu$ and $z-iu$ are each apart from zero
by $\Rea z$. The elementary Laplace integral for cosine gives
\[
 \int_0^\infty e^{-zs}\cos(ts)\dd s=\frac{z}{z^2+t^2}.
\]
It is verified by writing cosine as two exponentials and integrating
their finite antiderivatives. The sine integral on the right of the
next display has $t$-derivative equal to this cosine integral: on a
finite interval this is finite differentiation, and its derivative
tail is at most $e^{-aS}/a$. Comparing its value at $t=0$ with the
defining arctangent integral proves
\begin{equation}
 \operatorname{atan}(t/z)=
       \int_0^\infty e^{-zs}\frac{\sin(ts)}s\dd s,
 \qquad |\operatorname{atan}(t/z)|\le t/a
 \quad(\Rea z\ge a).
 \label{eq:atan-laplace}
\end{equation}
At $s=0$ use $\sin(ts)/s=t\operatorname{sinc}(ts)$. The estimate
$|\sin(ts)|\le ts$ for nonnegative real $s,t$ supplies both the
integrable majorant and the derivative justification.

Insert \eqref{eq:atan-laplace} into
$2\int_0^\infty\operatorname{atan}(t/z)/(e^{2\pi t}-1)\dd t$.
The resulting double integral has absolute majorant
\[
 \frac{2e^{-as}t}{e^{2\pi t}-1}.
\]
This is an explicit finite-sum argument: the $s>S$ tail is at most
$4e^{-aS}/a$, and the $t>T\ge1$ tail is at most
$4(T+1)e^{-T}/a$. The small strips $s<\eta$ and $t<\eta$
contribute at most $4\eta$ and $\eta/(\pi a)$ respectively.
These coarse bounds follow by splitting the remaining integral at $1$,
using $e^{2\pi t}-1\ge2\pi t$ below $1$ and
$(e^{2\pi t}-1)^{-1}\le2e^{-t}$ above it.
On the remaining finite rectangle the order of summation is a finite
rearrangement.

Expand the geometric denominator in the inner $t$ integral. Its $n$th
Laplace sine integral, after the regularized division by $s$, is
\[
 \int_0^\infty e^{-2\pi nt}\frac{\sin(st)}s\dd t
              =\frac1{s^2+4\pi^2n^2}.
\]
The tail after $n=N$ is at most $1/(4\pi^2N)$, uniformly in real
$s\ge0$, since $|\sin(st)/s|\le t$. Thus the inner sum is precisely
$k(s)/2$, proving
\begin{theorem}[Binet's second formula]
With the preceding right-half-plane branches,
\[
 \boxed{\displaystyle
 \mathcal L(z)=(z-\tfrac12)\operatorname{Log}z-z
 +\tfrac12\log(2\pi)
 +2\int_0^\infty
       \frac{\operatorname{atan}(t/z)}{e^{2\pi t}-1}\dd t.}
\]
All exchanges have been reduced to finite rectangles with the stated
tails. In particular no contour integral or branch continuation theorem
is needed.
\end{theorem}

For real $x>0$, the positive-kernel formula immediately proves the
additional consequence
\[
 (-1)^j R^{(j)}(x)=\int_0^\infty t^je^{-xt}k(t)\dd t>0
 \qquad(j\in\N).
\]
A positive rational margin follows by restricting the integral to
$[1,2]$ and using its first positive summand. This is complete
monotonicity of the Stirling remainder with an explicit constructive
meaning.

\section{Log-gamma integrals, Mellin transforms, and applications}
\label{sec:further}

\subsection{Raabe's formula and the unit-interval integral}

Let $J_s=\int_0^{\pi/2}\log(\sin t)\dd t$. These endpoint integrals
have effective tails: by $2t/\pi\le\sin t\le t$, an omitted interval
of length $\eta$ contributes at most
$\eta(1+|\log\eta|+\log2)$ in absolute value. Reflection about
$\pi/2$ gives the same bound for cosine. Substitution and the double
angle identity therefore give, first on truncated intervals,
\[
 2J_s=\int_0^{\pi/2}\log(\sin t\cos t)\dd t
     =-\tfrac\pi2\log2+\tfrac12\int_0^\pi\log(\sin u)\dd u
     =-\tfrac\pi2\log2+J_s.
\]
Hence $\int_0^1\log(\sin\pi t)\dd t=-\log2$.
The endpoint behavior
$\log\Gamma(t)=-\log t+\log\Gamma(1+t)$ is equally effective:
for $0<t\le1/2$, (C37)--(C38) give
$|\log\Gamma(1+t)|\le2$, so the omitted gamma-log integral is at
most $\eta(3+|\log\eta|)$. With $\eta=2^{-h}$, the choice
$h=4(p+10)$ makes each such tail less than $2^{-p}$.

Take real logarithms in the reflection identity and integrate on
$[\eta,1-\eta]$. Positivity licenses the logarithms; the preceding
tails allow the endpoints to be restored. We obtain
\begin{equation}
 \boxed{\displaystyle
 \int_0^1\log\Gamma(t)\dd t=\tfrac12\log(2\pi).}
 \label{eq:unit-loggamma}
\end{equation}
For $x>0$ put $I(x)=\int_x^{x+1}\log\Gamma(t)\dd t$.
The finite fundamental theorem and the recurrence give
$I'(x)=\log\Gamma(x+1)-\log\Gamma(x)=\log x$.
Thus $I(x)-x\log x+x$ is constant. Letting $x$ tend to zero with the
explicit endpoint estimates just proved identifies the constant in
\eqref{eq:unit-loggamma}. This proves Raabe's formula
\begin{equation}
 \boxed{\displaystyle
 \int_x^{x+1}\log\Gamma(t)\dd t
 =x\log x-x+\tfrac12\log(2\pi)\qquad(x>0).}
 \label{eq:raabe}
\end{equation}
The recurrence also gives
$\int_0^1\log\Gamma(1+t)\dd t=\tfrac12\log(2\pi)-1$.

\subsection{The gamma--zeta Mellin identity}

For $\Rea s\ge1+1/A$, define $\zeta(s)$ by its Dirichlet sums.
Their absolute tail is at most $A N^{-1/A}$, by integrating
$t^{-1-1/A}$ from $N$ onwards. In particular a cutoff
$N\ge(2^{p+4}A)^A$ gives a primitive recursive name. The integral
\[
 \int_0^\infty\frac{t^{s-1}}{e^t-1}\dd t
\]
has small-endpoint bound $A\eta^{1/A}$, since $e^t-1\ge t$,
and large-endpoint bound twice the gamma tail, since
$(e^t-1)^{-1}\le2e^{-t}$ for $t\ge1$.
For a finite $N$, integrate the exact finite geometric sum
$\sum_{n=1}^N e^{-nt}$ against $t^{s-1}$ and use the positive real
scaling substitution. The result is
$\Gamma(s)\sum_{n=1}^N n^{-s}$. The omitted integral is absolutely
bounded by
\[
 \Gamma(\Rea s)\sum_{n>N}n^{-\Rea s}
 \le (A+2^{R+1}R!)A N^{-1/A},\quad |s|\le R.
\]
This last comparison is established first for a finite sum of the
positive majorants and then by their stated geometric and integral
tails; it is not an appeal to a measure-theoretic interchange theorem.
It proves
\begin{equation}
 \boxed{\displaystyle
 \Gamma(s)\zeta(s)=\int_0^\infty\frac{t^{s-1}}{e^t-1}\dd t,
 \qquad\Rea s>1.}
 \label{eq:gamma-zeta-mellin}
\end{equation}
For $s=2$ this includes
$\int_0^\infty t/(e^t-1)\dd t=\pi^2/6$.
The zeta functional equation is a separate theorem and is not being
inferred from this half-plane identity.

\subsection{Wallis, Gaussian moments, and factorial interpolation}

At $z=1/2$ the sine product gives
\[
 \frac\pi2=\prod_{k=1}^{\infty}\frac{4k^2}{4k^2-1}.
\]
If $W_N$ is its finite product, then
$0\le\pi/2-W_N\le2/N$. To see the rate, the remaining logarithm is
at most $\sum_{k>N}2/(4k^2)\le1/(2N)$, and
$e^u-1\le2u$ for $0\le u\le1$. Use $W_N<\pi/2<2$.

The quadratic integral substitution and the half-integer values give,
for $c>0$ and $n\in\N$,
\[
 \int_{-\infty}^{\infty}t^{2n}e^{-ct^2}\dd t
   =\frac{(2n)!\sqrt\pi}{4^n n!c^{n+1/2}},\qquad
 \int_{-\infty}^{\infty}t^{2n+1}e^{-ct^2}\dd t=0.
\]
The odd integral is absolutely convergent, and its two finite symmetric
halves cancel exactly. In particular the Gaussian integral is
$\sqrt{\pi/c}$. Likewise the positive function
$p_x(t)=t^{x-1}e^{-t}/\Gamma(x)$ has integral one, and its finite
integer moments are
\[
 \int_0^\infty t^n p_x(t)\dd t=(x)_n.
\]
Thus its mean and variance are $x$ and $x$, by the cases $n=1,2$.
These statements are integral identities with the already proved
cutoffs; no additional probability-theoretic foundation is assumed.

Two equations mentioned in the reference article reduce to algebra.
For $z>1/2$,
$\Gamma(z-1/2)=\Gamma(z+1/2)$ holds exactly when $z=3/2$,
by recurrence and positive cancellation. For positive $z$ with
$z-1$ apart from the poles, the equality
$\Gamma(z-1)=\Gamma(z+1)$ is equivalent to $z(z-1)=1$.
Its unique positive solution is
$\phi=(1+\sqrt5)/2$. In fact the equation itself forces
$z-1=1/z>0$, and then positive square-root uniqueness solves
$(2z-1)^2=5$; there is no real-order trichotomy step.

Finally, normalization and recurrence alone do not characterize gamma.
For a nonzero rational $c$,
\[
 F_c(z)=\Gamma(z)\exp(c\sin(2\pi z))
\]
satisfies the same recurrence and factorial values. It differs from
gamma at $z=1/4$. Its analytic and arithmetic presentation follows from
the functions already constructed. This supplies explicit alternative
interpolations and explains why the convexity hypothesis in the
Bohr--Mollerup theorem carries mathematical content.
